\chapter*{Abstract}
\addcontentsline{toc}{chapter}{Abstract}
\thispagestyle{plain}

Network segmentation is widely promoted as one of the most cost-effective controls for limiting the lateral movement of attackers, yet many small and medium-sized organisations continue to operate flat networks because the security benefits are often described qualitatively rather than measured. This project addresses that gap by constructing a reproducible, containerised cyber range that recreates a realistic ten-host campus network in two configurations --- a flat single-subnet design and a six-zone segmented design enforced by an \texttt{iptables}-based gateway --- and by running an identical battery of automated tests against both. Four quantitative metrics are computed: firewall policy accuracy, blast-radius exposure, HTTP availability, and a derived reliability index. The flat architecture exhibits a 100\% blast-radius exposure (eight of eight peer hosts reachable from a compromised student workstation) and only 50\% policy accuracy, while the segmented architecture reduces exposure to 25\% and achieves 100\% policy accuracy. HTTP availability remains at 100\% in both topologies, with average response times of 1.03~ms (flat) and 1.08~ms (segmented), indicating that the security gains are obtained at a negligible performance cost (a difference of approximately 50~$\mu$s per request). The composite reliability index rises from 52.4 to 91.1 out of 100 once segmentation is applied. The work contributes (i) an open testbed that other students can clone in minutes on commodity hardware, (ii) a reusable measurement methodology aligned with NIST SP~800-41 and SP~800-207, and (iii) empirical evidence that supports the inclusion of segmentation in baseline security controls for resource-constrained environments. Limitations relating to the simplified threat model, the use of network-layer enforcement only, and the absence of host-level controls are discussed, together with directions for future work that would extend the testbed towards a full zero-trust architecture.

\vspace{1em}
\noindent\textbf{Keywords:} network segmentation; firewall; lateral movement; blast radius; cyber range; Docker; SEED Labs; zero trust; defence in depth; reproducible security testing.

